% Propietario conservado: /Users/ruben/Documents/New project/output/REVISION_ARGUMENTAL_APERTURAS_CIERRES_20260915/III/spanish_source/manuscrito/sections/03_hojas_y_orientacion.tex
\section[Retour de feuillet, incidence aréale--volumétrique et orientation constitutive]{Retour de feuillet, incidence aréale--volumétrique\\et orientation constitutive}
\label{iii:iii:sec:hojas}

La structure des feuillets intervient avant l’évaluation des grandeurs
électromagnétiques. Sa fonction consiste à distinguer la fermeture d’une projection,
la restitution d’une orientation et l’évolution du registre transporté.
Ces opérations agissent sur des objets liés, mais de types différents : une
trajectoire projetée peut s’être refermée tandis que son relèvement reste
sur le feuillet opposé ; l’orientation peut avoir été restituée tandis que la mémoire
conserve les deux tours effectués. La construction constitutive doit recevoir
cette information avec la structure qui la détermine.

\subsection{Retour projeté et restitution du revêtement orienté}

Soit $\ell$ une boucle de retour d’une extension survivante et soit
$\varepsilon:\langle\ell\rangle\to\{+1,-1\}$ le caractère de son
relèvement, avec $\varepsilon(\ell)=-1$. Le fibré d’intervalles associé
a pour présentation
\begin{equation}
 \mathcal M=([0,1]\times[-1,1])/{(0,u)\sim(1,-u)}.
 \label{iii:iii:eq:mobius}
\end{equation}
La donnée décisive est le signe du transport sur la boucle spécifiée.
L’identification des extrémités réalise géométriquement ce caractère.

\begin{proposition}[Deux ordres de clôture]
Un tour de la boucle ferme la base et inverse l’orientation de la fibre. Deux
tours restituent l’orientation. Dans la paramétrisation angulaire du lecteur de
bord, ces retours correspondent respectivement à la clôture aréale $2\pi$ et
à la clôture volumétrique $4\pi$ du revêtement orienté.
\end{proposition}
\begin{proof}
La relation d’équivalence de \eqref{iii:iii:eq:mobius} identifie la fibre terminale
à la fibre initiale au moyen de $u\mapsto-u$. La composition de deux transports est
$u\mapsto-(-u)=u$. Paramétrer un tour de la base par $2\pi$ attribue $4\pi$
à son carré. L’assertion concerne ce relèvement et cette lecture.
\end{proof}

La mémoire du TPK conserve en outre la trajectoire parcourue et les mises à jour du registre. Le carré du caractère de signe est l'identité, mais n'impose pas l'annulation de ces mises à jour. La représentation spinorielle, lorsqu'elle sera utilisée, devra transporter cette donnée au moyen de sa propre application.

\subsection{Descente de l’incidence depuis le calendrier tritique}

Dans le bloc primitif $(+1,-1,0)$, le TPK met à jour le mode arithmétique au moyen de
$+1:m\mapsto\Sigma$, $-1:m\mapsto\Pi$ et $0:m\mapsto m$. Son orbite cyclique
est $(\Sigma,\Pi,\Pi)$. Le lecteur de feuillets attribue $\Sigma$ au feuillet aréal et
$\Pi$ au feuillet volumétrique. Dans l’ordre des bases $(e_{\rm ar},e_{\rm vol})$,
l’incidence des arêtes est représentée par
\begin{equation}
 F_{\rm av}=\begin{pmatrix}0&1\\1&1\end{pmatrix}.
 \label{iii:iii:eq:fav}
\end{equation}
Les trois paires consécutives sont $\Sigma\to\Pi$, $\Pi\to\Pi$ et $\Pi\to\Sigma$, chacune de multiplicité un. Les calendriers de longueurs $9$, $27$ et $108$ contiennent $3$, $9$ et $36$ copies du bloc primitif. Leurs matrices d'incidence sont donc ces multiples de $F_{\rm av}$.

Cette construction conserve la différence entre fréquence chronologique et mesure
de l’espace des continuations. La première assigne des poids $(1/3,2/3)$ au cycle
des modes. La seconde se définit sur l’enveloppe symbolique minimale de mémoire
un qui admet exactement les couples précédents. Sa matrice d’adjacence est
$F_{\rm av}$ ; les dénombrements de trajectoires appartiennent à cette enveloppe.

\Needspace{10\baselineskip}
\begin{proposition}[Poids de retour des deux feuillets]
Soient $F_0=0$, $F_1=1$ et $F_{n+1}=F_n+F_{n-1}$. Pour $n\geq1$,
\begin{equation}
 F_{\rm av}^{n}=
 \begin{pmatrix}F_{n-1}&F_n\\F_n&F_{n+1}\end{pmatrix}.
 \label{iii:iii:eq:fav-powers}
\end{equation}
La mesure de Parry de cette enveloppe attribue aux feuillets des poids proportionnels à $(1,\varphi^2)$, et le rapport des décomptes diagonaux converge vers $\varphi^{-2}$.
\end{proposition}
\begin{proof}
La formule des puissances se prouve par induction en utilisant la récurrence.
Le carré de $F_{\rm av}$ est strictement positif. Son vecteur positif de
Perron est $(1,\varphi)$, de valeur propre $\varphi$, car
$\varphi^2=\varphi+1$. La matrice étant symétrique, les poids de Parry sont
proportionnels aux carrés des composantes de ce vecteur. L’autre racine
du polynôme caractéristique est $-\varphi^{-1}$ ; son module est strictement
inférieur à $\varphi$, ce qui donne la limite des quotients diagonaux.
\end{proof}

L’identification spectrale de cette incidence est postérieure à sa construction
à partir du calendrier. La coordonnée régionale $\varphi$ conserve sa généalogie
coinductive antérieure : la reconnaissance de la valeur propre ne sélectionne pas le calendrier.

\subsection{Marquage de la frontière et rapport surfacique--volumétrique}

Soit $\Pi_{\rm HMT}(X_\infty)$ la valeur du lecteur régional de clôture déjà construit. Avec les projecteurs diagonaux de feuillet, on définit
\begin{equation}
 D_\pi=\Pi_{\rm HMT}(X_\infty)P_{\rm ar}+P_{\rm vol},\qquad
 \Theta_n=
 \frac{\langle e_{\rm ar},D_\pi F_{\rm av}^{n}e_{\rm ar}\rangle}
 {\langle e_{\rm vol},F_{\rm av}^{n}e_{\rm vol}\rangle}.
 \label{iii:iii:eq:marked-return}
\end{equation}
La normalisation volumétrique est unitaire ; la clôture régionale marque une seule fois
le retour aréal. En appliquant \eqref{iii:iii:eq:fav-powers},
\begin{equation}
 \Theta_n=\Pi_{\rm HMT}(X_\infty)\frac{F_{n-1}}{F_{n+1}}
 \longrightarrow\frac{\pi}{\varphi^2}.
 \label{iii:iii:eq:av-limit}
\end{equation}
Le rapport de retour est déterminé par l'incidence, la mesure et le marquage. Le quotient $2\pi/4\pi$ décrit, en revanche, les deux ordres de clôture du revêtement. Les deux relations participent à la lecture des feuillets sans se substituer l'une à l'autre.

\subsection{Transport de l’orientation aux canaux constitutifs}

Dans une représentation à deux canaux, soit $\mathcal R$ une involution
autoadjointe et $\mathcal S$ un opérateur unitaire d’échange qui satisfait
$\mathcal S\mathcal R\mathcal S^{-1}=-\mathcal R$. Les coordonnées angulaires
déjà construites s’écrivent $x=A_{\rm rad}$ et $y=C^*_{\rm rad}$, avec $x>|y|$.
Alors
\begin{equation}
 B=xI+y\mathcal R,\qquad T=e^{-B},\qquad
 \mathcal ST(y)\mathcal S^{-1}=T(-y).
 \label{iii:iii:eq:channel-conjugation}
\end{equation}
En effet, la conjugaison transforme $B(y)$ en $B(-y)$ et commute avec sa série
exponentielle convergente. Les valeurs propres $q_\pm=e^{-x\mp y}$ sont échangées.
Pour fixer l’ordre des réponses, on utilise la chambre $x>y>0$, où
$0<q_+<q_-<1$. L’application $(x,y)\mapsto(x,-y)$ la transporte dans la
chambre opposée $x>-y>0$, où $0<q_-<q_+<1$. Toutes deux appartiennent au domaine
contractif $x>|y|$. À la frontière $y=0$, les deux canaux coïncident.
La conjugaison échange les chambres, elle ne conserve pas l’ordre strict de
leurs étiquettes ; elle agit sur les projecteurs au moyen de
$\mathcal SP_\pm\mathcal S^{-1}=P_\mp$, avec
$P_\pm=(I\pm\mathcal R)/2$.
Cette application concrète transmet l’orientation à l’opérateur constitutif
développé ci-après.
